% Source: https://artofproblemsolving.com/community/c3h1783551_interesting_question
% Stable identity: aops-topic-1783551-q01
Let $p$ be the polynomial. At each crossing point $p$ has a real root, so $p$ has at least $10$ distinct real roots and $\deg p\ge10$. The degree $10$ occurs, for example, for
\[
p(x)=(x-1)(x-2)\cdots(x-10),
\]
whose graph crosses the axis at $1,2,\dots,10$.

There is \emph{no maximum}: for every $k\ge0$, the polynomial $(x-1)(x-2)\cdots(x-10)\,(x^2+1)^k$ has degree $10+2k$ and the same sign as $(x-1)\cdots(x-10)$ everywhere, so its graph crosses the axis exactly at the same $10$ points.

On the other hand the degree must be \emph{even}. The graph crosses the axis exactly where $p$ changes sign, so between $x\to-\infty$ and $x\to+\infty$ the sign of $p$ changes exactly $10$ times, an even number. Hence $p(x)$ has the same sign for very large negative and very large positive $x$. For odd degree the leading term $cx^n$ has opposite signs at the two ends, so the degree cannot be odd.

Answer: the minimum degree is $\boxed{10}$, there is no maximum, and the possible degrees are exactly the even numbers $10,12,14,\dots$.
